%! AP Statistics — Unit 7, Lesson 7.2 — Warm-Up
\documentclass[10pt]{article}
\usepackage{apstats-article}
\usepackage{apstats-boxes}
\newcommand{\TallMath}[1]{$\displaystyle #1\rule[-1.4em]{0pt}{3.2em}$}

\begin{document}

\pageheader{Unit 7, Lesson 7.2}{Warm-Up}
\namedateperiod

\vspace{0.15in}

From Unit 6, the one-sample $z$-interval for a proportion has the form
$\hat p \pm z^* \sqrt{\hat p(1 - \hat p)/n}$.

\begin{enumerate}
  \item A random sample of 100 adults finds 64 support a new policy. Verify
        conditions and construct a 95\% CI for the true proportion.

        \begin{itemize}
          \item \textbf{Random:} \blank{5cm}
          \item \textbf{Large Counts:} $n\hat p = $ \blank{2cm}$ \ge 10$? \blank{1cm}
                \quad $n(1-\hat p) = $ \blank{2cm}$ \ge 10$? \blank{1cm}
          \item $\hat p = $ \blank{2cm} \quad $z^* = $ \blank{2cm}
          \item $95\%$ CI: \blank{8cm} $= ($ \blank{2cm} $,$ \blank{2cm} $)$
        \end{itemize}

  \item Now suppose we want to estimate a population \emph{mean} instead of a
        proportion. We know $\bar x = 42$, $s = 8$, $n = 36$, but $\sigma$ is unknown.

        \begin{enumerate}[label=\alph*.]
          \item Why can't we use a $z$-interval for this situation?
                \writelines{2}

          \item What distribution should we use instead, and what parameter(s)
                determine its shape?
                \writelines{2}
        \end{enumerate}
\end{enumerate}

\end{document}
